\section{Introduction}
\label{sec:introduction}

A bank's model-validation function runs tests on other people's models: the calibration of a default probability, the discrimination of a score, the exposure of a derivatives portfolio. The tests are code, and the code can be wrong. Supervisors ask the validator to benchmark and to confirm that the implementation reproduces the documented model \citep{fed2011sr117,ecb2025guide}, and the usual answer is a reconciliation: the validation engine is run on an input for which another source gives the number, and the two are compared. Four kinds of other source are in use. A paper prints a table, and the engine is fed the paper's inputs. A second computation of the same definition is written from scratch. A library publishes a closed form. A reference engine simulates the same model. Each comparison ends in an outcome that the engine agrees or does not, and the outcome is usually all that is reported.

An outcome is worth what the comparison behind it is worth, and three things decide that. The first is what the reference shares with the engine: a reference that calls the same library, or the same helper, cannot show a defect in it, whatever the numbers say. The second is the rule: when the reference is a simulation, the difference of the two numbers is random, and a rule that accepts it when it is small compared with the sampling error has a probability of rejecting an engine that is right and a probability of accepting one that is wrong. These are the size and the power of the rule, and they are seldom measured. The third is what is told about the gaps: a gap that is listed with its size in a unit the reader understands and with a tested explanation is information, and one that is summarised as a failure or dropped is not.

This work proposes a reconciliation protocol that answers the three, and measures it on a real engine. The engine is \texttt{riskvalidation} 1.1.0, the validation engine of the Contraste product for the validation of a bank's risk models. Its estimators and tests are reconciled with five papers of credit risk, with two recomputations of its tail probabilities that call none of its SciPy functions, with QuantLib's Hull--White closed forms, with the \texttt{transitionMatrix} library, and with the Open Source Risk Engine for the exposure of a swap book, with and without collateral.

The contributions are the following.
\begin{enumerate}
\item A protocol with one rule per class of reference (printed answer, independent route, closed form, Monte Carlo reference engine), a common report of the gap in a unit the reader can interpret with the explanation tested, an equivalence margin and a detectable gap beside every pass, and a list of what is not claimed (Section~\ref{sec:proposed-method}).
\item Six propositions that state what each rule can establish: the digit rule and its power, the blindness of a comparison that reuses the engine's function, the size of a simultaneous rule over many dates, the miscalibration of the standard error of antithetic pairs, the detectable gap, and the detection threshold of a relative gate (Section~\ref{sec:theory}).
\item A reconciliation of \nPrintedN\ printed numbers in which \nPrintedAgree\ agree to the digit and each of the \nMissTotal\ gaps is listed, \nMissTested\ of them explained by a tested rounding of the paper's own inputs; independent routes that agree with the engine's $p$-values to $\nGridJeffMax$ and that, unlike a reuse of the engine's call, flag a planted error from $10^{-9}$ upward; closed forms that agree with QuantLib on \nClosedWithin\ of \nClosedCases\ cases (Section~\ref{sec:results}).
\item The first measurement of the size and power of the acceptance rule that a Monte Carlo reconciliation uses: the product's rule rejects two runs of one model in \nSizeProdNetting\,\% of the repetitions, because the antithetic pairs make its errors too large and the thirty-five comparisons make its bound too small, and a refinement with pair-aware errors and a family-wise bound has a size of \nSizeRefineNetting\,\% and a power of \nPowerRefineSigmaTwo\,\% against a $2\,\%$ volatility error, where the product's rule has \nPowerProdSigmaTwo\,\% (Sections~\ref{sec:results-calibration} and~\ref{sec:results-power}).
\item The reported gaps of the collateralised variants against ORE, reproduced on the product's draws and shown to be properties of the two engines on \nReplCsaN\ new seeds, with the part that a zero minimum transfer amount explains (Section~\ref{sec:results-seeds}).
\end{enumerate}

The main result is that the reconciliations that exist hold, with qualifications that are reported where they arise, and that the rule which the Monte Carlo reconciliation uses has a size of \nSizeProdNetting\,\% and a power of \nPowerProdSigmaTwo\,\% against a $2\,\%$ volatility error, which are numbers a validation function can use. The remainder of the manuscript places the protocol among related work (Section~\ref{sec:related-work}), formulates the problem (Section~\ref{sec:problem}), states the protocol and its theory (Sections~\ref{sec:proposed-method} and~\ref{sec:theory}), describes the experiments and their results (Sections~\ref{sec:experimental-design} to~\ref{sec:ablation}), and discusses limits and reproducibility (Sections~\ref{sec:discussion} to~\ref{sec:reproducibility}).
